\section{Background and Problem Formulation}
\label{sec:background}

\subsection{MoE Inference and Expert Offloading}
\label{sec:background_moe}

A Mixture-of-Experts (MoE) layer replaces the dense feed-forward network
with a set of experts and a router. For each input token, the router selects
a small subset of experts and combines their outputs according to the routing
weights~\cite{shazeer2017outrageously,fedus2022switch,
jiang2024mixtral,dai2024deepseekmoe}. We refer to an expert's model-defined
identity as its \emph{logical expert ID}. For an MoE invocation at step $t$,
we denote by $\mathcal{A}_t$ the set of logical experts selected by the
router across all participating tokens. Because routing is data dependent,
$\mathcal{A}_t$ can change across inference steps. The effect is particularly
pronounced during prefill, where many prompt tokens are processed together
and their routed experts may form a much larger active set than in a typical
decoding step~\cite{fang2025klotski}.

When the complete expert pool does not fit in accelerator HBM, expert
offloading keeps the expert weights in host memory and maintains only a
bounded resident working set on the accelerator
~\cite{xue2024moeinfinity,hwang2024pregated,cao2025moelightning,
fang2025klotski}. Let $\mathcal{H}_t$ denote the set of experts resident in
HBM at step $t$, with $|\mathcal{H}_t| \leq C$ for an HBM capacity of $C$
experts. A routed expert in $\mathcal{A}_t \cap \mathcal{H}_t$ can be used
directly, whereas an expert in
$\mathcal{A}_t \setminus \mathcal{H}_t$ incurs an offloading miss and must
be transferred from host memory. If the HBM working set is already full,
loading the missing expert requires replacing an existing resident expert.

Importantly, offloading makes expert \emph{placement} a runtime-managed
state. Although the logical identity of expert $e$ is fixed by the model,
the physical storage containing its weights need not be fixed across
inference steps. We denote the physical location of a resident expert $e$
at step $t$ by $\pi_t(e)$. Expert loading and replacement can therefore
change both the resident set $\mathcal{H}_t$ and the placement mapping
$\pi_t$ as routing evolves. This distinction between stable logical expert
identity and dynamic physical placement is central to the interaction
between expert offloading and graph replay discussed next.

\subsection{Graph Replay and Replay-Stable Execution State}
\label{sec:background_graph}

Graph replay reduces repeated host-side scheduling and launch overhead by
capturing a recurring accelerator execution sequence and reusing it across
subsequent invocations. Once captured, however, the replayed computation is
tied to the execution structure and graph-visible storage bindings recorded
during capture. We refer to these reusable properties collectively as the
\emph{replay-stable execution state}. In this paper, this state includes the
captured operation topology and the physical storage referenced by captured
computation.

Replay stability does not imply that all runtime data must remain unchanged.
The contents of graph-visible buffers may evolve between replays as long as
the captured computation continues to access the same storage. For example,
FlashInfer accommodates request-dependent scheduling under CUDA Graph by
placing scheduler metadata in preallocated device workspaces and updating
their contents between invocations, while preserving the addresses referenced
by captured kernels~\cite{ye2025flashinfer}. This distinction separates
\emph{dynamic values} from \emph{dynamic bindings}: the former can be written
into replay-stable storage, whereas changing the physical storage referenced
by captured computation changes the execution state on which replay depends.

We therefore model a captured region by its operation topology
$\mathcal{O}$ and graph-visible physical bindings $\mathcal{B}$. Across
replays, runtime values may evolve from $V_t$ to $V_{t+1}$, while
$\mathcal{O}$ and $\mathcal{B}$ remain reusable:
\[
    (\mathcal{O}, \mathcal{B}, V_t)
    \;\longrightarrow\;
    (\mathcal{O}, \mathcal{B}, V_{t+1}).
\]
The key requirement is thus not immutable data, but a stable execution
interface through which changing runtime data can be consumed. Conventional expert offloading violates this separation because expert replacement can change the physical weight binding itself rather than merely
the contents stored behind a stable binding.

\subsection{Dynamic Expert Placement Meets Graph Replay}
\label{sec:background_conflict}

The interaction between expert offloading and graph replay becomes problematic
when captured computation directly references the physical storage of expert
weights. Consider a logical expert $e$ that is resident at physical location
$\pi_t(e)$ during step $t$. If $e$ is later evicted and subsequently reloaded,
the runtime may place it at a different location $\pi_{t+1}(e)$. Thus,
although the logical expert remains unchanged, its physical binding may evolve
across inference steps:
\[
    \pi_t(e) \neq \pi_{t+1}(e).
\]

This behavior conflicts with the replay model in
Section~\ref{sec:background_graph}. If captured expert computation directly
binds logical expert $e$ to its current physical weight storage, then changing
$\pi_t(e)$ also changes the graph-visible binding $\mathcal{B}$ on which the
captured execution depends. In other words, conventional offloading couples
two forms of state that have different stability requirements: logical expert
placement must evolve with routing and HBM residency, whereas graph-visible
physical bindings must remain reusable across replay.

Several straightforward alternatives avoid one side of this conflict only by
sacrificing another system objective. Keeping all experts permanently resident
would stabilize their physical bindings but defeats offloading under a bounded
HBM budget. Falling back to eager execution accommodates arbitrary placement
changes but forfeits the host-overhead savings of graph replay. Rebuilding or
repairing captured execution after placement changes similarly reintroduces
runtime intervention into the repeated inference path. The desired execution
model must therefore allow expert identity and residency to evolve while
preventing those changes from propagating into the physical bindings observed
by captured computation.